 \setcounter{subsection}{0}

\section{CONCLUSIONES}

\subsection{Análisis de Impacto del Proyecto}

En esta sección se analiza el impacto del trabajo en dos planos. Primero se describe lo que aporta a la ingeniería de sistemas y luego su efecto en los contextos global, económico, ambiental y social.

\subsubsection{Impacto en la ingeniería de sistemas}

Además del sistema construido, el trabajo deja tres prácticas de ingeniería que se pueden
trasladar a otros proyectos, y cada una tiene un límite. Primero, la arquitectura se eligió
comparando alternativas con el método ATAM \citep{bass_software_2021}: un monolito modular,
microservicios de grano fino y la variante de grano grueso adoptada. La matriz ponderada del
documento ATAM da 3,91 sobre 5 a la variante elegida, frente a 3,44 y 2,82; el acta del
segundo objetivo había registrado 4,27 en una evaluación anterior, con otro conjunto de
alternativas. La decisión se apoyó en impulsores comprobables (recuperabilidad asimétrica,
frontera de cumplimiento y perfil de recurso heterogéneo) y no en la escalabilidad invocada
al principio, que el cálculo de carga descartó por órdenes de magnitud.

Segundo, el sistema se auditó más allá de sus requisitos funcionales. Las dos auditorías del
segundo objetivo revisaron el código de los servicios anteriores. La de autorización encontró
22 de 152 puntos de entrada sin una comparación completa de ámbito, y la de datos personales,
catorce hallazgos, entre ellos políticas de seguridad a nivel de fila que fallaban abiertas
cuando faltaba el contexto de usuario. El sistema construido corrigió una parte: ya no devuelve
los valores rechazados en los errores, no copia el correo del acudiente en los avisos y no
guarda el texto de la lista oficial. El ADR 0015 había decidido copiar las claves de ámbito
en cada tabla protegida y escribir contra ellas las políticas de fila; el sistema construido no
tiene esas políticas y controla el ámbito en la aplicación. Los
anexos de auditoría de autorización y de datos personales registran la situación de cada
hallazgo, y la Sección~\ref{sec:defectos} reúne los defectos que siguen abiertos. Tercero, las
fronteras entre servicios pasaron de convención documental a funciones de aptitud ejecutables
\citep{ford_building_2017}: doce pruebas contrastan el código con el manifiesto y otras
veintiocho comprueban que esas pruebas detectan infracciones plantadas. La integración continua
las ejecuta, aunque 11 de sus 12 ejecuciones registradas hasta el final de septiembre
terminaron en fallo y su estado en la versión evaluada no se comprobó. Las tres prácticas sirven a un proyecto que deba sostener ante terceros que su
diseño hace lo que declara, siempre que el seguimiento de los hallazgos forme parte del
método.

\subsubsection{Impacto en contextos globales, económicos, ambientales y sociales}

El problema que dio origen a Tukit es social antes que tecnológico. La encuesta del primer
objetivo específico (N=65 respuestas recibidas, n=60 analizables) describe una compra
costosa para los acudientes: el 65\,\% le dedica cuatro horas o más y el 68,3\,\% compra en
tres o más lugares. Las listas son, además, heterogéneas (1.451 formas distintas de nombrar
los mismos artículos sobre 3.608 renglones reunidos de 17 colegios), y cada hogar las
interpreta por su cuenta. Si el sistema reduce ese esfuerzo, el efecto sería económico, al
facilitar la comparación de precios y sustitutos que el proceso manual no ofrece, y de tiempo
para las familias en las semanas de mayor presión del calendario escolar. Ese efecto y la carga
cognitiva son lo que mide el OE5 (Sección~\ref{sec:oe5})\pendienteEvidencia{resultados SUS, NPS y carga cognitiva del OE5 para el impacto social}. El modelo de negocio adoptado
\citep{osterwalder_business_2010}, B2B2C, distribuye ese valor sin imponer un costo al colegio:
la institución no paga por participar, y el ingreso proviene de la compra del acudiente. En el
contexto colombiano, donde la capacidad de pago de las instituciones y de sus hogares es
heterogénea, una solución que dependiera del presupuesto del colegio limitaría su adopción a
las instituciones con mayor capacidad financiera. En el plano ambiental, el trabajo no midió un efecto directo. La disponibilidad que optimiza el
cuarto objetivo es la fracción de líneas del kit que Tukit puede surtir, un criterio comercial. Que
reduzca compras redundantes es una hipótesis cuya medición exige operar con familias, y la economía
circular de libros y útiles se plantea como trabajo futuro (Sección~\ref{sec:trabajo-futuro}). El alcance social del hallazgo tiene un límite: la encuesta se
concentró en un solo colegio del sector Salitre, con sobrerrepresentación del estrato 4
(68,3\,\%) y ausencia de los estratos 1 y 6. Por eso, generalizar la insatisfacción declarada a
la población escolar colombiana excede lo que esta muestra permite afirmar.

\subsection{Conclusiones y Trabajo Futuro}

En esta sección se evalúa el cumplimiento de cada objetivo específico y su razón, y luego se
presenta el trabajo futuro. La Tabla~\ref{tab:cumplimiento} resume, por objetivo, la evidencia, la
sección donde se presenta y el estado, y los párrafos siguientes explican por qué se asignó ese
estado.

\begin{table}[H]
\centering\footnotesize
\setlength{\tabcolsep}{3pt}
\caption{Cumplimiento por objetivo específico.}
\label{tab:cumplimiento}
\begin{tabularx}{\textwidth}{@{}p{0.9cm}Xp{2.9cm}p{2.3cm}@{}}
\toprule
OE & Evidencia & Sección & Estado \\
\midrule
OE1 & Corpus de 17 colegios, 3.608 renglones, 1.451 formas de nombrar y 99 conceptos canónicos; encuesta (n=60) & Cap.~III y IV & Cumplido \\
OE2 & Comparación ATAM de tres alternativas, árbol de utilidad de 17 escenarios, vistas 4+1 y 38 ADR & Cap.~V & Cumplido \\
OE3 & F1 micro de 0,948 en lectura (CGS-10); resolución con $\tau$ calibrado; NL2C & Sec.~\ref{sec:percepcion}, \ref{sec:optimizacion-cal} y~\ref{sec:nl2c} & Cumplido, con las amenazas de la Sec.~\ref{sec:amenazas} \\
OE4 & Optimizador de cuatro objetivos, igualdad con el óptimo exacto en 174 listas, MAUT, hipervolumen y aplicación web con pago de pruebas & Sec.~\ref{sec:optimizacion} y~\ref{sec:optimizacion-cal} & Cumplido\pendienteEvidencia{estado final de OE4 tras MAUT, hipervolumen y NL2C} \\
OE5 & Pruebas de unidad, integración, sistema y aceptación con usuarios (SUS, NPS, RTLX) & Sec.~\ref{sec:control-calidad} y~\ref{sec:oe5} & Cumplido\pendienteEvidencia{estado final de OE5 segun C1, C2 y C3 del plan de analisis} \\
\bottomrule
\end{tabularx}
\end{table}

El OE1 se cumplió porque el corpus permitió definir 99 conceptos canónicos, de los cuales 25
cubren cerca del 80\,\% de los renglones observados: el dominio es heterogéneo en la superficie,
pero tiene una estructura tratable. El OE2 se cumplió porque las decisiones de arquitectura se
justificaron con ATAM \citep{bass_software_2021}, vistas 4+1 \citep{kruchten_41_1995} y la norma
ISO/IEC/IEEE 42010 \citep{iso_42010_2022}; el diseño previó siete desplegables y el sistema
construido tiene seis, agrupados por razón de cambio, sobre cinco bases de datos (ADR 0013).

El OE3 se cumplió con un cambio de diseño entre los prototipos y el sistema construido. En
percepción, el modelo multimodal superó al reconocimiento con reglas (F1 micro de 0,948 frente a
0,298), con la salvedad de que el oro se construyó desde su salida. En resolución, el sistema
construido sustituyó la fusión recíproca de rangos y el reordenador por una sola fase que suma el
parecido léxico y el coseno, y compensa con la revisión de cada lista nueva antes de publicarla. El
OE4 se cumplió porque el optimizador enumera el frente exacto hasta 4.096 combinaciones y recurre a
NSGA-II \citep{deb_fast_2002,blank_pymoo_2020} por encima, con cuatro objetivos (costo, fidelidad
de marca, disponibilidad y personaje), y presenta los kits Ahorrador y Premium siempre, Fan si el
acudiente pidió un personaje y A tu Medida si eligió el punto medio de gasto o un tope. El personaje
es un hallazgo del desarrollo: la encuesta midió esa preferencia y los tres criterios de la VFP no la
recogían. El OE5 se evaluó con el método de la VFP y el plan de análisis fijado antes de la primera
sesión, y su estado depende de los criterios C1, C2 y C3 de ese plan.

\subsubsection{Trabajo futuro}
\label{sec:trabajo-futuro}

El trabajo futuro se organiza en seis líneas, alineadas con la sección de trabajo futuro de la
propuesta (VFP v4). Cada una quedó fuera de este trabajo por una razón explícita, que se indica a
continuación.
\begin{enumerate}
\item \textbf{Alianza con tiendas de uniformes e integración de libros.} Alianza con las tres
tiendas de uniformes y flujo con Buscalibre y las editoriales de textos escolares. Quedó fuera
porque este trabajo se centró en un solo negocio, el de útiles y papelería.
\item \textbf{Aplicaciones nativas y modo sin conexión.} El acudiente usa una aplicación web
progresiva; las aplicaciones nativas y el modo sin conexión amplían su alcance, pero no cambian el
algoritmo que se evaluó.
\item \textbf{Endurecimiento productivo.} Autorización uniforme en las rutas, rol de acceso por base
de datos, depuración del nombre del menor en el documento cargado, supresión de datos, medición de
la cobertura en la integración continua y cobro real con la pasarela. Es la condición para operar
con familias reales y se plantea antes del despliegue.
\item \textbf{Optimización de rutas.} Agrupar la entrega entre tiendas, que no afecta la
construcción del kit.
\item \textbf{Economía circular de libros.} Reventa e intercambio de textos entre hogares, que
requiere una base de usuarios activa.
\item \textbf{Calibración con más anotadores y más colegios.} Repetir la calibración del umbral y
de los pesos con un conjunto anotado por dos personas y con listas de más colegios, y comparar
NSGA-II con otros optimizadores de referencia sobre el mismo banco.
\end{enumerate}
